\documentclass[11pt,a4paper]{article}
\usepackage[utf8]{inputenc}
\usepackage[T1]{fontenc}
\usepackage[french]{babel}
\usepackage{lmodern}
\usepackage{microtype}
\usepackage{geometry}
\usepackage{hyperref}
\hypersetup{colorlinks=true, urlcolor=blue, linkcolor=black, citecolor=black}
\usepackage{amsmath, amssymb}
\usepackage{graphicx}
\usepackage{enumitem}
\usepackage{titlesec}
\usepackage{fancyhdr}
\setlength{\headheight}{14pt}
\usepackage{tabularx}
\usepackage{xcolor}
\usepackage{array}
\usepackage{setspace}
\usepackage{listings}
\usepackage{float}
\usepackage{colortbl}
\usepackage{booktabs}
\usepackage{tcolorbox}
\usepackage{fontawesome5}
\usepackage{tikz}
\usetikzlibrary{arrows.meta, positioning, fit, backgrounds}

% ==============================
% MISE EN PAGE ET COULEURS
% ==============================
\geometry{left=2.2cm, right=2.2cm, top=2.2cm, bottom=2.2cm}
\setstretch{1.05}
\setlist[itemize]{itemsep=2pt, topsep=2pt, parsep=0pt}

\definecolor{codeblue}{rgb}{0.13,0.40,0.72}
\definecolor{codegreen}{rgb}{0.10,0.55,0.10}
\definecolor{codeorange}{rgb}{0.80,0.40,0.00}
\definecolor{backgray}{rgb}{0.96,0.96,0.98}
\definecolor{blueTN}{RGB}{0,75,150}
\definecolor{lightblue}{RGB}{220,235,250}
\definecolor{lightgreen}{RGB}{225,245,225}
\definecolor{lightyellow}{RGB}{255,245,220}
\definecolor{lightred}{RGB}{252,232,232}
\definecolor{purpleTN}{RGB}{95,55,140}

\tcbuselibrary{skins,breakable}

\newtcolorbox{infocadre}[1][]{
  colback=lightblue!60, colframe=blueTN, fonttitle=\bfseries,
  title={\faInfoCircle\ #1}, arc=4pt, boxrule=0.8pt, left=6pt, right=6pt, top=4pt, bottom=4pt, breakable
}
\newtcolorbox{stepcadre}[1]{
  colback=lightgreen!50, colframe=codegreen, fonttitle=\bfseries,
  title={\faTerminal\ #1}, arc=3pt, boxrule=0.7pt,
  left=6pt, right=6pt, top=3pt, bottom=3pt, breakable
}
\newtcolorbox{probcadre}[1]{
  colback=lightred!60, colframe=red!60!black, fonttitle=\bfseries,
  title={\faExclamationTriangle\ #1}, arc=3pt, boxrule=0.7pt,
  left=6pt, right=6pt, top=3pt, bottom=3pt, breakable
}
\newtcolorbox{choixcadre}[1]{
  colback=lightyellow!70, colframe=codeorange, fonttitle=\bfseries,
  title={\faLightbulb\ #1}, arc=3pt, boxrule=0.7pt,
  left=6pt, right=6pt, top=3pt, bottom=3pt, breakable
}

\lstdefinelanguage{BashShell}{
  keywords={python3, cd, source, docker, compose, up, ps, ssh, rsync, sudo, systemctl, journalctl, iptables, grep, host, cat, mosquitto_sub, echo},
  keywordstyle=\color{codeblue}\bfseries, sensitive=true,
  comment=[l]{\#}, commentstyle=\color{codegreen}\itshape,
  stringstyle=\color{codeorange}, morestring=[b]",
}
\lstset{
  basicstyle=\scriptsize\ttfamily, keywordstyle=\color{codeblue}\bfseries,
  commentstyle=\color{codegreen}\itshape, stringstyle=\color{codeorange},
  numbers=none, backgroundcolor=\color{backgray}, frame=single,
  rulecolor=\color{blueTN!30}, breaklines=true, breakatwhitespace=true,
  tabsize=2, showstringspaces=false, xleftmargin=6pt, xrightmargin=4pt,
  aboveskip=4pt, belowskip=3pt,
}

% ==============================
% EN-TÊTES ET TITRES
% ==============================
\pagestyle{fancy}
\fancyhf{}
\fancyhead[L]{\small\bfseries\color{blueTN} AI Smart Parental Gateway}
\fancyhead[R]{\small Compte Rendu \& Guide Démo}
\fancyfoot[C]{\textcolor{blueTN}{\thepage}}
\renewcommand{\headrulewidth}{0.6pt}
\renewcommand{\footrulewidth}{0.3pt}

\titleformat{\section}{\Large\bfseries\color{blueTN}}{\thesection}{0.6em}{}[\color{blueTN}\titlerule]
\titlespacing*{\section}{0pt}{12pt}{5pt}
\titleformat{\subsection}{\large\bfseries\color{codeblue}}{\thesubsection}{0.6em}{}

\begin{document}

% ==============================
% PAGE DE GARDE
% ==============================
\thispagestyle{empty}
\begin{center}
\vspace*{1.5cm}
\begin{tcolorbox}[
  colback=lightblue!35, colframe=blueTN, arc=6pt, boxrule=1.2pt,
  width=0.9\textwidth, halign=center, valign=center,
  top=16pt, bottom=16pt, left=10pt, right=10pt
]
{\fontsize{20pt}{24pt}\selectfont\textbf{\faNetworkWired\ GuardianAI Gateway}}\\[10pt]
{\large\textit{Rapport de Projet \& Guide de Démonstration Officiel}}
\end{tcolorbox}
\vspace{1cm}
{\Large\textbf{Projet de Fin d'Année}}\\[0.5cm]
{\large Réalisé par : \textbf{Rayen FEHRI}}\\
{\large Encadrante : \textbf{Mme Zaineb KAMMOUN}}\\
\vspace{0.5cm}
{\large Institut Supérieur d'Informatique}\\
\vspace{1.5cm}

\begin{choixcadre}{Aperçu du Système}
\textbf{Architecture :} Raspberry Pi (hostapd + dnsmasq + iptables + systemd) + PostgreSQL + MQTT + Grafana + Intelligence Artificielle (Isolation Forest \& Mistral AI).\\
Ce document fusionne le compte rendu d'architecture et le script étape par étape de la démonstration technique.
\end{choixcadre}

\end{center}
\newpage

\tableofcontents
\newpage

% ==============================
% SECTION 1 : THEORIE (Compte Rendu 2)
% ==============================
\section{Partie I : Architecture et Moteur de Règles}

L'objectif de cette passerelle est d'appliquer dynamiquement des politiques de restriction réseau. L'orchestrateur \texttt{main\_rules.py} synchronise l'état désiré (base de données) avec l'état réel (pare-feu et résolveur DNS).

\subsection{Modèle de Données (PostgreSQL)}
Une table dédiée centralise l'ensemble des règles de contrôle.

\begin{center}
\small
\begin{tabularx}{\textwidth}{|l|l|X|}
\hline
\rowcolor{lightblue}\multicolumn{3}{|c|}{\textbf{Table \texttt{rules} — Politiques de contrôle parental}} \\ \hline
\rowcolor{lightblue!40}\textbf{Colonne} & \textbf{Type} & \textbf{Rôle} \\ \hline
\texttt{id} & SERIAL (PK) & Identifiant unique de la règle \\ \hline
\texttt{mac} & VARCHAR & Adresse MAC ciblée (NULL = tous) \\ \hline
\texttt{rule\_type} & VARCHAR & \texttt{block\_device}, \texttt{schedule}, \texttt{block\_category} \\ \hline
\texttt{active} & BOOLEAN & État d'activation (\texttt{TRUE}/\texttt{FALSE}) \\ \hline
\end{tabularx}
\end{center}

\subsection{Blocage Réseau (iptables)}
L'interdiction repose sur le blocage des trames au niveau MAC.
\begin{itemize}
    \item \textbf{Ciblage MAC :} Règle visant la chaîne \texttt{FORWARD}.
    \item \textbf{Idempotence :} Vérification via \texttt{iptables -C} avant ajout.
    \item \textbf{Plages horaires :} Gestion intelligente du chevauchement de minuit.
\end{itemize}

\subsection{Filtrage DNS Catégoriel (dnsmasq)}
\begin{itemize}
    \item \textbf{Mappage JSON :} Association de domaines à des catégories (Réseaux sociaux, Jeux, etc.).
    \item \textbf{Poisoning local :} Redirection via \texttt{address=/domaine.com/0.0.0.0}.
    \item \textbf{Redémarrage intelligent :} Comparaison MD5 avant de recharger le service DNS.
\end{itemize}

\section{Partie II : Intelligence Artificielle}

L'architecture IA est \textbf{hybride} et se divise en deux volets :

\subsection{Détection d'Anomalies (Règles + Machine Learning)}
L'orchestrateur IA (\texttt{main\_ai.py}) analyse l'historique DNS :
\begin{itemize}
    \item \textbf{Détecteurs Déterministes :} Alertes immédiates sur les accès nocturnes ou les catégories sensibles (ex: Paris sportifs).
    \item \textbf{Isolation Forest :} Modèle non supervisé (via \texttt{scikit-learn}) qui apprend le comportement moyen de l'appareil (volume, horaires, répartition) et signale les écarts inattendus.
\end{itemize}

\subsection{Génération de Rapports (Mistral AI)}
L'API LLM de \textbf{Mistral AI} traduit ces données techniques pour le parent :
\begin{enumerate}
    \item \textbf{Agrégation} des requêtes et anomalies sur N jours.
    \item \textbf{Génération} d'un rapport avec profilage du temps d'écran.
    \item \textbf{Proactivité :} L'IA émet des suggestions directes prêtes à l'emploi (ex: \textit{"Bloquer les réseaux sociaux de 22h à 7h pour l'adolescent"}).
\end{enumerate}

\newpage

% ==============================
% SECTION 2 : PRATIQUE (Demo Commands)
% ==============================
\section{Partie III : Guide de Démonstration Pas-à-Pas}

\begin{center}
\small
\begin{tabularx}{\textwidth}{|l|l|X|}
\hline
\rowcolor{blueTN}\multicolumn{3}{|c|}{\textcolor{white}{\textbf{Tableau de bord de l'infrastructure}}} \\ \hline
\rowcolor{lightblue!40}\textbf{Service} & \textbf{Adresse / Accès} & \textbf{Identifiants} \\ \hline
Raspberry Pi (SSH) & \texttt{ssh pi@raspberrypi.local} & \texttt{pi} / \texttt{ray+ray000} \\ \hline
Grafana & \texttt{http://localhost:3000} & \texttt{admin} / \texttt{guardian123} \\ \hline
PostgreSQL & \texttt{localhost:5433} & \texttt{gateway\_admin} / \texttt{changeme123} \\ \hline
MQTT Broker & \texttt{localhost:1883} & Anonyme \\ \hline
WiFi du Pi & SSID : \texttt{GuardianGateway} & -- \\ \hline
\end{tabularx}
\end{center}

\subsection{Étape 0 — Démarrage (PC Fedora)}
\begin{stepcadre}{Lancement des conteneurs}
\begin{lstlisting}[language=BashShell]
cd ~/9raya/Projects/gateway_project
docker compose up -d
docker ps --format "table {{.Names}}\t{{.Status}}\t{{.Ports}}"
\end{lstlisting}
\end{stepcadre}

\subsection{Étape 1 — Résilience (Raspberry Pi)}
Le routeur intelligent s'auto-configure au démarrage via \texttt{systemd}.
\begin{stepcadre}{Vérification des services autonomes}
\begin{lstlisting}[language=BashShell]
ssh pi@raspberrypi.local
sudo systemctl status guardian-network guardian-discovery guardian-rules --no-pager
\end{lstlisting}
\end{stepcadre}

\subsection{Étape 2 — Blocages (Total et Horaire)}
Via l'ordinateur, nous modifions directement la base de données. Le Raspberry Pi détecte le changement et applique la règle au réseau instantanément.
\begin{stepcadre}{Blocage MAC et Horaire (iptables)}
\begin{lstlisting}[language=BashShell]
cd ~/9raya/Projects/gateway_project
source venv/bin/activate
MAC="a2:50:d5:8f:0b:04"

# Blocage total immédiat :
python3 -c "from db import add_rule; add_rule('$MAC', 'block_device')"

# Blocage planifié (horaire) :
python3 -c "from db import add_rule; add_rule('$MAC', 'schedule', '22:00', '07:00')"
\end{lstlisting}
\end{stepcadre}

\subsection{Étape 3 — Blocage DNS Catégoriel}
\begin{stepcadre}{Blocage des Réseaux Sociaux}
\begin{lstlisting}[language=BashShell]
python3 -c "from db import add_rule; add_rule(None, 'block_category', 'reseaux_sociaux')"
\end{lstlisting}
\end{stepcadre}

\begin{infocadre}[Vérification DNS]
Un test local sur le Pi via \texttt{host tiktok.com 127.0.0.1} renvoie désormais l'IP \texttt{0.0.0.0}, prouvant l'efficacité du blocage, alors que \texttt{google.com} reste accessible.
\end{infocadre}

\newpage
\subsection{Étape 4 — Démonstration de l'IA (Injection et Analyse)}
Afin de valider le Machine Learning et la génération de rapport LLM, un trafic familial réaliste est simulé sur 7 jours.
\begin{stepcadre}{Pipeline IA Complet}
\begin{lstlisting}[language=BashShell]
# 1. Injection des données de test
python3 inject_test_data.py

# 2. Lancement des détecteurs (SQL + Isolation Forest)
python3 main_ai.py

# 3. Lancement de la génération du rapport via Mistral AI
python3 main_ai.py --report --days 7
\end{lstlisting}
\end{stepcadre}

\begin{itemize}
    \item \textbf{Résultat attendu 1 :} Le modèle \textit{Isolation Forest} et les règles SQL détectent l'usage nocturne du smartphone de l'adolescent et la tentative d'accès à un site de casino.
    \item \textbf{Résultat attendu 2 :} L'appel à Mistral AI génère un rapport en langage naturel avec des suggestions proactives, telles que : \textit{"Suggestion : Bloquer l'accès aux réseaux sociaux entre 22h et 7h pour l'appareil Galaxy-A54."}
\end{itemize}

\subsection{Étape 5 — Télémétrie MQTT et Grafana}
\begin{stepcadre}{Écoute en temps réel des événements réseau}
\begin{lstlisting}[language=BashShell]
docker exec gateway_mosquitto mosquitto_sub -t "guardian/#" -v
\end{lstlisting}
\end{stepcadre}
En parallèle, le Dashboard Grafana accessible sur \texttt{http://localhost:3000} synthétise toutes ces données (Top sites visités, règles en vigueur, appareils en ligne).

% ==============================
% DÉPANNAGE
% ==============================
\section{Dépannage Rapide (Troubleshooting)}

\begin{center}
\small
\begin{tabularx}{\textwidth}{|l|l|X|}
\hline
\rowcolor{lightred!50}\textbf{Symptôme} & \textbf{Vérification} & \textbf{Solution} \\ \hline
Services Pi inactifs & \texttt{systemctl status guardian-*} & \texttt{sudo systemctl start ...} \\ \hline
Pas d'Internet client & \texttt{iptables -t nat -L POSTROUTING} & Vérifier NAT (MASQUERADE) \\ \hline
Blocage inefficace & \texttt{iptables -L FORWARD -n} & Vérifier priorité règle DROP \\ \hline
Erreur Mistral AI & \texttt{cat .env} & Vérifier la variable \texttt{MISTRAL\_API\_KEY} \\ \hline
\end{tabularx}
\end{center}

\vspace{1cm}
\begin{center}
\textbf{\large --- Fin de la démonstration ---}
\end{center}

\end{document}
